% angr false-positive root causes vs HOFPR proofs (7 SAFE samples)
\begin{table}[t]
\centering
\small
\begin{tabular}{lp{5.2cm}p{4.2cm}}
\toprule
Sample & angr FP root cause & HOFPR proof \\
\midrule
  safe_01_guarded_strcpy & 2 & strcpy 调用即报, 不分析源长度与守卫 & src_bound \\
  safe_02_clamp_alloc & 1 & memcpy 目标为堆对象即报, 不解析 clamp 后的分配大小 & const_cmp \\
  safe_03_literal_strcpy & 1 & strcpy 调用即报, 不识别字面量源 & src_bound \\
  safe_04_loop_bounded & 1 & 循环内写堆即报, 不关联循环边界与分配大小 & loop_shared_var \\
  safe_05_aligned_slack & 1 & 按请求大小 17B 判定, 忽略 ptmalloc2 对齐余量 (usable=24) & const_cmp \\
  safe_06_shared_var & 1 & memcpy 大小参数为变量即报, 不建立分配-写入共享变量关联 & shared_var \\
  safe_07_guard_var_alloc & 1 & 分配大小为变量乘积即报, 不追踪 clamp 守卫上界与仿射关系 & affine_dom \\
  safe_08_affine_double & 1 & memcpy 大小变量即报, 不建立分配-写入仿射支配 (2n>=n) & affine_dom \\
  safe_09_affine_pad & 1 & memcpy 大小变量即报, 不建立分配-写入仿射支配 (n+64>=n) & affine_dom \\
\bottomrule
\end{tabular}
\caption{Why generic binary tools report every heap write as an
overflow, and how HOFPR eliminates each false positive with a
formal proof. angr flags 7/7 safe samples; HOFPR proves all safe.}
\label{tab:fp_cause}
\end{table}
